\part{微分几何}
在德国形而上学的经典论述中, 对世界进行抽象的重点并不在于文本, 符号, 结构, 秩序, 而在于一个先于秩序的本源世界, 康德将这个本源世界称为"物自体(Ding an sich)". 在康德那里, 物自体作用于主体, 使主体感知到外部刺激并得到未经加工的原始感官素材, 随后通过感性、知性、理性3种认识论处理原始素材, 其中感性能力通过时间与空间两种先天直观形式把握现象并得到经验直观, 知性能力通过12种先天认识范畴将这些感性直观进一步综合得到经验知识, 理性能力则将经验知识整合、统一形成一般规律. 康德将那些原始感官素材全体构成的集合称为"杂多\footnote{蓝公武, 纯粹理性批判, 1957}(Mannigfaltigkeit)", 黎曼经由高斯及费希特的学生赫尔巴特(Herbart)了解了康德的哲学体系, 并在1851年将"杂多"用到了自己的博士论文《单复变函数的一般理论基础》\footnote{\href{https://www.emis.de/classics/Riemann/Grund.pdf}{Bernhard Riemann, "Grundlagen für eine allgemeine Theorie der Functionen einer veränderlichen complexen Grösse", 1851}}中\sidenote{一些文本误认为黎曼首次将"杂多"引入数学是在1854年的演讲《关于几何基础中的假设》\footnote{\href{https://www.emis.de/classics/Riemann/Geom.pdf}{Bernhard Riemann, Ueber die Hypothesen, welche der Geometrie zu Grunde liegen, 1854}}中, 但实际上首次引入是在1851年的博士论文中.}. 1900年胡塞尔在《逻辑研究》中拓展了"杂多", 将这个概念与黎曼的工作融合, 形成了现象学的理论基础. 洋务运动过后, 陈省身的师兄江泽涵参与翻译了欧洲的数学著作, 他将黎曼论文中的"杂多"译为中文"流形", 这一译法沿用至今. 一种说法是, 江泽涵从文天祥的《正气歌》中吸收了灵感, 天地有正气, 杂然赋流形. 下则为河岳, 上则为日星. 於人曰浩然, 沛乎塞苍冥. 不论这一传言是否符合史实, 这个源于德国古典哲学与的词语都在事实上深刻影响了微分几何学的精神气质.

\input{main/differential_geometry/coordinates.tex}
\newpage
\input{main/differential_geometry/connection.tex}
\newpage
